\documentclass[11pt]{article}
\usepackage{fontspec}
\usepackage[french]{babel}
\usepackage{amsmath}
\usepackage{amssymb}
\usepackage{amsthm}
\usepackage{mathrsfs}
\usepackage{tikz-cd}
\usepackage{enumitem}

\newcommand{\printedpage}[1]{\marginpar{\footnotesize p.~#1}}
\newcommand{\msp}[1]{\marginpar{\footnotesize [#1]}}
\newcommand{\cM}{\mathcal{M}}
\newcommand{\Flu}{\underline{\mathrm{Fl}}}
\newcommand{\Fl}{\mathrm{Fl}}
\newcommand{\fU}{\mathfrak{U}}
\newcommand{\fV}{\mathfrak{V}}

\begin{document}

\printedpage{4}
% head: A. Grothendieck : Dérivateurs, Ch. XV, Édition des Universités de Montpellier II et Paris VII
% foot: Version 5 novembre 2010

Ainsi, si on considère
\[
\left\{
\begin{array}{rl}
\Flu(\cM) & \text{comme objet de CAT}/\cM \text{ par le foncteur } f \longmapsto \text{but}\, f\\[1ex]
\text{où} & \text{CAT est la catégorie des catégories \textit{éléments de} } \fV\text{\textit{,} univers}\\
 & \text{tel que } \fU \in \fV\text{, d'où Cat } \overset{\text{(déf)}}{=} \mathrm{Cat}^{\fU},\ \cM,\ \Flu(\cM) \in \text{CAT} ,
\end{array}
\right.
\]
on a un endofoncteur $p_1$ de $\Flu(\cM)$ sur $\cM$, et un homomorphisme fonctoriel
\[
i_1 \ : \ \mathrm{id}_{\Flu(\cM)/\cM} \longrightarrow p_1 \ , \qquad \textit{i.e.} \quad
\begin{tikzcd}[column sep=small]
X \arrow[rr, "f"] \arrow[dr, "i_1(f)"'] & & Y \\
 & \underbrace{X_1(f)}_{=\,\text{source}\, p_1(f)} \arrow[ur, "p_1(f)"'] &
\end{tikzcd} \ ,
\tag*{(1.13)}
\]
et la description de la suite ordinale des endofoncteurs $p_\alpha$ de $\Flu\,\cM/\cM$ résulte de cette seule donnée, compte tenu que les petites limites filtrantes existent dans les fibres de $\Flu\,\cM \xrightarrow{\text{but}} \cM$.

\msp{page 5}

\[
\begin{tikzcd}[column sep=normal]
X \arrow[r, "i_1(f)"] \arrow[dr, "f"'] & X_1(f) \arrow[r, "i_1(p_1(f))"] \arrow[d, "p_1(f)"] & X_2(f) \arrow[r] \arrow[dl, "p_2(f)"] & \cdots \arrow[r] & X_\alpha(f) \arrow[r, "i_1(p_\alpha(f))"] \arrow[dlll, "p_\alpha(f)"] & X_{\alpha+1}(f) \arrow[r] \arrow[dllll, "p_{\alpha+1}(f) = p_1(p_\alpha(f))"] & \cdots \\
 & Y & & & & &
\end{tikzcd}
\tag*{(1.14)}
\]

Donc, pour achever de définir la suite ordinale des $X_\alpha(f)$ (munie des $i_\alpha(f)$, $p_\alpha(f)$), ou des $p_\alpha(f)$ (munie des $i_\alpha(f)$), il reste à décrire $X_1(f)$, $i_1(f)$, $p_1(f)$.

La propriété essentielle de ce triple dont nous aurons besoin est la suivante :

\begin{itemize}[leftmargin=5em, labelwidth=4em, labelsep=1em, align=left]
\item[(1.15)] \textbf{Fact}$(C_0)$ (propriété de factorisation par rapport à $C_0$).

Pour tout diagramme commutatif
\[
D : \qquad
\begin{tikzcd}
A \arrow[r] \arrow[d, "i"'] & X \arrow[d, "f"] \\
B \arrow[r] & Y
\end{tikzcd}
\qquad \text{avec } i \in C_0 \ ,
\]
d'où un diagramme traits pleins
\[
\begin{tikzcd}
A \arrow[r] \arrow[d, "i"'] & X \arrow[d, "i_D"] \arrow[dr, dotted, "i_1(f)"] \arrow[dd, dotted, bend right=60, "f"'] & \\
B \arrow[r] & X_D \overset{\text{déf}}{=} X \amalg_A B \arrow[r, dotted, "u_D"] \arrow[d, "p_D"] & X_1(f) \arrow[dl, dotted, "p_1(f)"] \\
 & Y \ , &
\end{tikzcd}
\]
\printedpage{5}
% head: A. Grothendieck : Dérivateurs, Ch. XV, Édition des Universités de Montpellier II et Paris VII
% foot: Version 5 novembre 2010
il existe un homomorphisme de la factorisation $f = p_D i_D$ de $f$ dans la factorisation $f = p_1(f)i_1(f)$.
\end{itemize}

Pour formuler cette condition, j'ai utilisé l'existence des sommes amalgamées $X_D$, \textit{i.e.} la condition sur $C_0$ :

\msp{page 6}

\begin{itemize}[leftmargin=5em, labelwidth=4em, labelsep=1em, align=left]
\item[(1.16)] Les flèches $i \in C_0$ sont coquarrables, \textit{i.e.} pour un diagramme
\[
\begin{tikzcd}
A \arrow[r, "u"] \arrow[d, "i"'] & X \\
B &
\end{tikzcd}
\]
avec $i \in C_0$, la somme amalgamée existe.
\end{itemize}

Mais on peut formuler 1.15 aussi ainsi, indépendamment de l'hypothèse (1.16) :

\begin{itemize}[leftmargin=5em, labelwidth=4em, labelsep=1em, align=left]
\item[(1.17)] Pour tout carré commutatif
\[
\begin{tikzcd}
A \arrow[r, "u"] \arrow[d, "i"'] & X \arrow[d, "f"] \\
B \arrow[r, "v"] & Y
\end{tikzcd}
\qquad ,
\]
avec $i \in C_0$, il existe un carré commutatif traits pleins
\[
\begin{tikzcd}
A \arrow[r, "u"] \arrow[d, "i"'] & X \arrow[d, "i_1(f)"] \\
B \arrow[r, "v_1"] \arrow[dr, dotted, "v"] & X_1(f) \arrow[d, dotted, "p_1(f)"] \\
 & Y
\end{tikzcd}
\]
tel que $v = p_1(f)v_1$.
\end{itemize}

Ceci a pour conséquence le

\textbf{Corollaire 1} (de 1.17). \textit{Soit} $f \in \Fl(\cM)$\textit{, et} $\alpha$ \textit{un ordinal} $(\in \fU)$. \textit{Pour tout carré commutatif}
\[
\begin{tikzcd}
A \arrow[r, "u_\alpha"] \arrow[d, "i"'] & X_\alpha(f) \arrow[d, "p_\alpha(f)"] \\
B \arrow[r, "v"] & Y
\end{tikzcd}
\]
\printedpage{6}
% head: A. Grothendieck : Dérivateurs, Ch. XV, Édition des Universités de Montpellier II et Paris VII
% foot: Version 5 novembre 2010
\textit{il existe un carré commutatif}
\[
\begin{tikzcd}
A \arrow[r, "u_\alpha"] \arrow[d, "i"'] & X_\alpha(f) \arrow[d, "\text{trans.}"] \arrow[dd, dotted, bend left=60, "p_\alpha(f)"] \\
B \arrow[r, "v_1"] \arrow[dr, dotted, "v"] & X_{\alpha+1}(f) \arrow[d, dotted, "p_{\alpha+1}(f)"] \\
 & Y
\end{tikzcd}
\]
\textit{tel que}
\[
v \ = \ p_{\alpha+1}(f)\ v_1 \ .
\]

\msp{page 7}

\textbf{Corollaire 2.} \textit{Soit $\pi$ un cardinal} $(\in \fU)$ \textit{tel que pour toute flèche $i : A \longrightarrow B$ dans $C_0$, la source $A$ soit $\pi$-accessible. Soit $\alpha$ un ordinal grand devant $\pi$} (i.e. \textit{tel que l'ensemble ordonné $I_\alpha$ des ordinaux $\alpha' < \alpha$ soit grand devant $\pi$}) -- \textit{p. ex. $\alpha$ le premier ordinal dont le cardinal soit $> \pi$. Alors pour toute flèche $f \in \Fl(\cM)$, on a $p_\alpha(f) \in (C_0)_*$.}

\textsc{Démonstration}. Cela signifie que pour tout carré commutatif
\[
\begin{tikzcd}
A \arrow[r, "u"] \arrow[d, "i"'] & X_\alpha(f) \arrow[d, "p_\alpha(f)"] \\
B \arrow[r, "v"] \arrow[ur, dotted, "w"] & Y \ ,
\end{tikzcd}
\tag*{$(*)$}
\]
avec $i \in C_0$, la flèche pointillée $w$ (avec $wi = u$, $p_\alpha(f)w = v$) existe. Or
\[
X_\alpha(f) = \varinjlim_{\alpha' \in I_\alpha} X_{\alpha'}(f) \ ,
\]
donc $A$ étant $\pi$-accessible et $I_\alpha$ grand devant $\pi$, $u$ provient d'une flèche $u_{\alpha'} : A \longrightarrow X_{\alpha'}(f)$, et (quitte à remplacer $\alpha'$ par $\alpha'' > \alpha'$ et $u_{\alpha'}$ par le composé $u_{\alpha''} : A \longrightarrow X_{\alpha'}(f) \longrightarrow X_{\alpha''}(f)$) on peut supposer $p_{\alpha'}(f)u_{\alpha'} = vi$, \textit{i.e.} que $(*)$ provient d'un carré commutatif
\[
\begin{tikzcd}
A \arrow[r, "u_{\alpha'}"] \arrow[d, "i"'] \arrow[rr, dotted, bend left=30, "u"] & X_{\alpha'}(f) \arrow[r, dotted, "\text{trans.}"'] \arrow[d, "p_{\alpha'}(f)"'] & X_\alpha(f) \arrow[dl, dotted, "p_\alpha(f)"] \\
B \arrow[r, "v"] & Y \ , &
\end{tikzcd}
\]

\printedpage{7}
% head: A. Grothendieck : Dérivateurs, Ch. XV, Édition des Universités de Montpellier II et Paris VII
% foot: Version 5 novembre 2010
et appliquant le corollaire 1, on trouve
\[
\begin{tikzcd}
A \arrow[r, "u_{\alpha'}"] \arrow[d, "i"'] & X_{\alpha'}(f) \arrow[d, "\text{trans.}"] & \\
B \arrow[r, "v_1"] \arrow[dr, dotted, "v"] & X_{\alpha'+1}(f) \arrow[r, dotted, "\text{trans.}"] \arrow[d, "p_{\alpha+1}(f)"] & X_\alpha(f) \arrow[dl, dotted, "p_\alpha(f)"] \\ % [sic] p_{\alpha+1}(f) -> p_{\alpha'+1}(f), the arrow starts at X_{\alpha'+1}(f); ?? uncertain: subscript of the label, tightly set -- \alpha+1 or \alpha'+1
 & Y &
\end{tikzcd}
\]
avec $v = p_{\alpha+1}(f)v_1$, et on prendra % [sic] p_{\alpha+1}(f)v_1 -> p_{\alpha'+1}(f)v_1, corollaire 1 is applied to \alpha' (arrow from X_{\alpha'+1}(f))
\[
w \ = \ \mathrm{trans}_{\alpha,\alpha'+1} \circ v_1 \ .
\]

\msp{page 8}

Mais on veut une factorisation de $f$, flèche quelconque dans $\cM$, en
\[
f \ = \ p\ i \ , \qquad \text{avec} \ \ p \in (C_0)_* \ \ \textit{et} \ \ i \in \widetilde{C_0} \ ,
\]
donc il faut encore énoncer une propriété, assurant que les $i_\alpha(f)$ sont dans $\widetilde{C_0}$. On aura même quelque chose de plus précis. Introduisons

\begin{itemize}[leftmargin=10em, labelwidth=9em, labelsep=1em, align=left]
\item[(1.18)\hfill $C'_0 \ = $] plus petit sous-ensemble de $\Fl(\cM)$ contenant $C_0$, stable par cochangement de base (chaque fois qu'il existe).
\item[\hfill $C''_0 \ = $] ensemble des flèches qui sont soit dans $C'_0$, soit des sommes amalgamées d'une famille \textit{non vide} de telles flèches $A \xrightarrow{i_\alpha} B_\alpha$ de même source $A$, $i_\alpha \in C'_0$ ; la flèche $A \longrightarrow B = \displaystyle\coprod_{\alpha}{}_{(A)} B_\alpha$ est dans $C''_0$.
\end{itemize}

(La condition 1.16 implique que les éléments de $C'_0$ sont coquarrables, et cela suffit, avec la condition de stabilité de $\cM$ par $\varinjlim$ filtrantes, à assurer que les sommes amalgamées infinies envisagées existent dans $\cM$.)

\textbf{N.B.} $C''_0$ est \textit{égal} au plus petit ensemble de flèches contenant $C_0$ et stable par cochangement de base et par petites sommes amalgamées. \texttt{[Si les petites sommes existent dans $\cM$, alors $C''_0$ est aussi l'ensemble des flèches obtenues par cochangement de base à partir de sommes de flèches de $C_0$.]}

\setcounter{footnote}{1}% the edition numbers these footnotes 2 and 3
\begin{itemize}[leftmargin=10em, labelwidth=9em, labelsep=1em, align=left]
\item[(1.19)\hfill $C'''_0 \ = $] plus petit sous-ensemble de $\Fl(\cM)$ contenant $C''_0$, et stable de plus par limites inductives ordinales strictes, \textit{i.e.} pour un système inductif ordinal $(A_\alpha)_{\alpha<\alpha_0}$, tel que pour $\alpha$ ordinal limite on ait $A_\alpha \xleftarrow{\ \sim\ } \varinjlim\limits_{\alpha' \in I_\alpha} A_{\alpha'}$, si les morphismes de transition $i_{\beta,\alpha} : A_\alpha \longrightarrow A_\beta$ sont dans $C'''_0$, de même les morphismes canoniques
\[
A_\alpha \xrightarrow{\ i_\alpha\ } \varinjlim_{\beta \in I_{\alpha_0}} A_\beta \ . \quad \footnotemark
\]
\end{itemize}
\footnotetext{$C_0 \subset C'_0 \subset C''_0 \subset C'''_0$}% footnote 2: mark printed in the display, text printed at the foot of p. 8

\printedpage{8}
% head: A. Grothendieck : Dérivateurs, Ch. XV, Édition des Universités de Montpellier II et Paris VII
% foot: Version 5 novembre 2010
\msp{page 9}

La proposition 2, p. 4 de XIII assure que l'on a
\[
C'''_0 \ \subset \ \widetilde{C_0} \ .
\]
Ceci posé, voici la deuxième propriété essentielle de la factorisation $f = p_1(f)i_1(f)$ (la première étant 1.17) :

\begin{itemize}[leftmargin=5em, labelwidth=4em, labelsep=1em, align=left]
\item[(1.20)] \qquad\qquad Pour toute flèche $f \in \Fl(\cM)$, on a $i_1(f) \in C''_0$.
\end{itemize}

Il en résulte aussitôt, par récurrence ordinale :

\begin{itemize}[leftmargin=5em, labelwidth=4em, labelsep=1em, align=left]
\item[(1.21)] \qquad Pour toute flèche $f \in \Fl(\cM)$ et tout ordinal $\alpha \in \fU$, on a $i_\alpha(f) \in C'''_0$, a fortiori $i_\alpha(f) \in \widetilde{C_0}$.
\end{itemize}

Donc si $\alpha$ est grand devant $\pi$ (\textit{cf.} cor. 2, page 7), la factorisation $f = p_\alpha(f)i_\alpha(f)$ satisfait
\[
\left\{
\begin{array}{l}
f = p_\alpha(f)i_\alpha(f)\text{, avec } i_\alpha(f) \in C'''_0 \subset \widetilde{C_0} = \Phi \text{ \texttt{[à ne pas}}\\
\text{\texttt{confondre avec }}\Phi_\alpha\text{\texttt{]}, } p_\alpha(f) \in C_{0*} = \Psi \text{, lorsque } \alpha \text{ est}\\
\text{grand devant } \pi\text{, les sources des } i \in C_0 \ \pi\text{-accessibles.}
\end{array}
\right.
\tag*{(1.22)}
\]

Il reste à construire le foncteur $f \longmapsto X_1(f)$, avec $i_1(f)$, $p_1(f)$, satisfaisant les conditions (1.17) et (1.20), sous l'hypothèse (1.16) que les $i \in C_0$ sont coquarrables. On considère alors, pour $f \in \Fl(\cM)$ fixé, l'ensemble $\Delta(f)$ (ou $\Delta(C_0, f)$) de \textit{tous} les diagrammes $D$ de la forme (1.15), avec $i \in C_0$, et on pose
\[
X_1(f) \ = \ \coprod_{D \in \Delta(f)}{}_{(X)} X_D
\tag*{(1.23)}
\]
(somme amalgamée sous $X$, \textit{i.e.} somme indexée
\msp{page 10}
par $\Delta(f)$ dans $X\backslash\cM$), le morphisme $X \xrightarrow{i_1(f)} X_1(f)$ étant le morphisme costructural, et le morphisme $X_1(f) \xrightarrow{p_1(f)} Y$ étant celui ayant pour composantes les $p_D$. En somme, on fait “au plus économique” (et au plus naturel) pour assurer la validité de (1.17), sous la forme équivalente (1.15). D'autre part, par définition de $C'_0$, $C''_0$, les $i_D$ sont dans $C'_0$, et $i_1(f)$ est dans $C''_0$, ce qui est (1.20)\footnote{Il faudrait quand même expliciter la dépendance fonctorielle de $X_1(f)$ par rapport à $f$\,!!!}.

Pour résumer :

\textbf{Théorème 1.} \textit{Soient $\cM$ une $\fU$-catégorie stable par $\fU$-limites inductives filtrantes, $C_0 \subset \Fl(\cM)$ un petit ensemble de flèches de $\cM$. On suppose les flèches $i \in C_0$ coquarrables. Cela assure la définition d'un endofoncteur $f \longmapsto p_1(f)$ de $\Flu(\cM)$, et d'un homomorphisme fonctoriel $i_1 : \mathrm{id}_{\Flu(\cM)} \longrightarrow p_1$,} i.e. \textit{$i_1(f) : f \longrightarrow p_1(f)$, induisant l'identité sur les objets buts de $f$ et de $p_1(f)$,}
\[
\begin{tikzcd}[column sep=small]
X \arrow[rr, "i_1(f)"] \arrow[dr, "f"'] & & X_1(f) \arrow[dl, "p_1(f)"] \\
 & Y \ , &
\end{tikzcd}
\]

\end{document}
